\documentclass[11pt,a4paper]{article}

%====================================================
% Packages
%====================================================

\usepackage{mathptmx}
\usepackage[left=0.75in,right=0.75in,top=0.7in,bottom=0.7in]{geometry}
\usepackage{titlesec}
\usepackage[hidelinks]{hyperref}
\usepackage{enumitem}

%====================================================
% Formatting
%====================================================

\setlength{\parindent}{0pt}
\setlength{\parskip}{0.35em}

\titleformat{\section}
{\large\bfseries}
{}{0em}{}[\titlerule]

\begin{document}

%====================================================
% Header
%====================================================

\begin{center}

{\LARGE \textbf{Younsoo Park}}

\vspace{0.3em}

\href{mailto:danielpark0605@gmail.com}{danielpark0605@gmail.com} \
\href{mailto:yqp5187@psu.edu}{yqp5187@psu.edu} \
+1 (814) 810-1707 \
GitHub: \href{https://github.com/danielsoo}{danielsoo}

\end{center}

%====================================================
% Education
%====================================================

\section*{Education}

\textbf{The Pennsylvania State University} \hfill University Park, PA

Bachelor of Science in Computer Science \
Minor in Mathematics \
Expected Graduation: May 2027 \
Dean's List

%====================================================
% Research Interests
%====================================================

\section*{Research Interests}

Systems Security \\
Network Security \\
Distributed Systems \\
LLM Systems \\
AI Agent Memory \\
Federated Learning \\
TinyML for IoT Security \\
Adversarial Machine Learning \\
Edge AI Systems

%====================================================
% Research Experience
%====================================================

\section*{Research Experience}

\textbf{Undergraduate Research Assistant \& Team Lead} \hfill Oct 2025 -- Present

\textit{Reliable Federated TinyML Deployment for IoT Security (WIP, LCTES '26)} \

Penn State University and Elizabethtown College \
Supervisors: Dr. Suman Saha and Dr. Peilong Li

\vspace{0.3em}

Research focus:
Federated and adversarially robust TinyML systems for IoT security and resource-constrained edge-device deployment.

\vspace{0.3em}

Key contributions:

\begin{itemize}[leftmargin=1.5em]

\item Designed a federated learning pipeline using FedAvgM, focal loss, and cosine learning-rate scheduling on CIC-IDS2017.

\item Improved attack recall from 46.7\% to 93.85\% and F1-score from 84.1\% to 89.32\%.

\item Built a TinyML compression pipeline including BatchNorm folding, structured pruning, knowledge distillation, and INT8 QAT/PTQ.

\item Achieved 12.28$\times$ model compression and 74.5\% inference latency reduction for ESP32 deployment.

\item Evaluated adversarial robustness using PGD attacks (\(\varepsilon\)-sweep 0.01--0.2) across 48 compression configurations.

\end{itemize}

%====================================================
% Publications
%====================================================

\section*{Publications and Manuscripts}

\textbf{Reliable Federated TinyML Deployment for IoT Security} \
Work in Progress \
Target Venue: ACM LCTES 2026



%====================================================

% Independent Study

%====================================================

\section*{Independent Study}

\textbf{AI Agent Memory Systems} \hfill Jul 2026 -- Present

Advisor: Prof. Kiwan Maeng

\begin{itemize}[leftmargin=1.5em]

\item Surveying production AI agent memory systems including Mem0 and A-MEM.

\item Investigating memory extraction, retrieval, safety, and lifecycle management for LLM-based agents.

\item Applying memory architectures to production recommendation systems through a Memory-first Personalization framework, shipped as Wowness at Levit Inc.

\end{itemize}

%====================================================
% Industry and Technical Experience
%====================================================

\section*{Industry and Technical Experience}

\textbf{Levit Inc. (Shopport AI)} \hfill Seoul, South Korea

Associate Problem Solver \hfill June 2026 -- Aug 2026

\begin{itemize}[leftmargin=1.5em]

\item Worked directly with the CEO to define product direction, prioritize engineering initiatives, and translate customer problems into production AI systems for an LLM-powered shopping assistant.

\item Built a production evaluation platform analyzing \textbf{49,807 conversations} and \textbf{80,124 query-level routing decisions}, achieving \textbf{97.7\% evaluation coverage} through versioned routing audit and historical backfill pipelines that reprocessed 47,231 conversations spanning 100+ days of production history.

\item Developed end-to-end experimentation infrastructure across frontend, backend, workers, and admin systems, improving recommendation completion from \textbf{62.8\%} to \textbf{75.0\% (+12.2 percentage points)}.

\item Evolved Shopport's routing stack through three architectural generations -- LLM classification, an embedding-based semantic fast path ($\sim$8$\times$ faster for confident cases), and a hybrid classifier for ambiguous queries -- reaching \textbf{86\% V3 routing accuracy} and rescuing \textbf{93\%} of eligible router errors.

\item Designed and shipped Wowness, a Mem0-backed cross-category discovery recommender grounded in the live catalog, as sole engineer across a 919-line NestJS engine, React UI, and admin review lab -- extending the Memory-first Personalization architecture from proposal to production.

\item Diagnosed production failures at the system boundary, including a split MongoDB connection that held the V3 feature switch at zero executions, and corrected systematic analytics false positives, reducing one queue-level count from 2,005 to 811.

\item Contributed \textbf{730 GitHub commits} and \textbf{143 merged pull requests}, adding over \textbf{114,000 lines} of TypeScript/TSX during the engagement.

\end{itemize}

\vspace{0.5em}

\textbf{SIGNUM (Nittany AI Alliance -- PIT-UN Funded)} \hfill Aug 2024 -- Nov 2025

Co-founder of a \$3,750 PIT-UN funded AI platform for caregiver decision support.

\begin{itemize}[leftmargin=1.5em]

\item Built multi-source ETL pipelines integrating CMS, NPPES, and Google Places datasets into a DuckDB warehouse.

\item Developed hospital-rating prediction models using Markov Transition analysis with confidence intervals.

\item Designed and deployed domain-specific RAG systems using FAISS and AWS Bedrock, reducing caregiver research time from 50+ hours per week to near-instant answers.

\end{itemize}

\vspace{0.5em}

\textbf{Atom Tech Solutions LTD} \hfill Berkeley, CA

Software Engineering Intern \hfill May 2022 -- Aug 2022

\begin{itemize}[leftmargin=1.5em]

\item Built secure authentication systems using SQL/JWT authentication and encrypted credential storage.

\item Implemented automated credential generation and authenticated onboarding workflows.

\end{itemize}
%====================================================
% Technical Projects
%====================================================

\section*{Technical Projects}

\textbf{GladIEEEators -- Penn State IEEE Battle Bots (1st Place)} \hfill Apr 2026

\begin{itemize}[leftmargin=1.5em]

\item Designed Shot \& Chaser, a coordinated two-robot combat system, for the IEEE Student Chapter Battle Bots Competition, winning 1st place.

\item Engineered a belt-driven AR500 vertical-spinner weapon system (334 J at 18,500 RPM, 200+ mph tip speed), validated via Fusion FEA above 4,000 N with a 2.09 minimum safety factor.

\item Designed mixed-material armor (6061 aluminum, UHMW, TPU) and integrated drivetrain, radio control, and protected LiPo electronics across both robots within a \$976 total build budget.

\end{itemize}

\vspace{0.5em}

\textbf{XIIO -- LinkedIn for Filmmakers} \hfill Apr 2026 -- Present

\begin{itemize}[leftmargin=1.5em]

\item Built a creator networking platform featuring portfolios, video streaming, collaboration workflows, and school-based discovery using Next.js, Firebase, and Cloudflare Stream.

\item Designed real-time messaging, creator profiles, media uploads, and production portfolio management.

\end{itemize}

\vspace{0.5em}

\textbf{ASME @ Penn State Website} \hfill Aug 2024 -- Present

\begin{itemize}[leftmargin=1.5em]

\item Developed the official ASME Penn State website using React, TypeScript, Vite, and Firebase.

\item Designed a multi-tier permission and approval system for organization management.

\end{itemize}

\vspace{0.5em}

\textbf{Hangukgwan (Family Business Web App)} \hfill 2024 -- Present

\begin{itemize}[leftmargin=1.5em]

\item Developed a multilingual full-stack restaurant web application using React, TypeScript, Node.js/Express, and MongoDB Atlas.

\item Implemented JWT authentication, Google Maps API integration, and multilingual i18n support.

\end{itemize}

%====================================================
% Teaching Experience
%====================================================

\section*{Teaching Experience}

\textbf{Lead Teaching Assistant, Math 230 (Multivariable Calculus)} \
Penn State University \hfill Aug 2025 -- Present

Responsibilities include assisting instruction, supporting students during office hours, and helping coordinate course activities.

%====================================================
% Leadership and Service
%====================================================

\section*{Leadership and Service}

\textbf{THON Technology Committee} \
Technology Captain (2027) \hfill Apr 2026 -- Present

\vspace{0.4em}

\textbf{Penn State Advanced Vehicle Team (AVT)} \
Project Manager \hfill Fall 2026 -- Present

\vspace{0.4em}

\textbf{Association for Computing Machinery (ACM)} \
Web \& Design Director \hfill Apr 2026 -- Present

\vspace{0.4em}

\textbf{Republic of Korea Air Force} \
Squad Leader \hfill Sept 2022 -- June 2024

%====================================================
% Relevant Coursework
%====================================================

\section*{Relevant Coursework}

\textbf{Computer Science}

\begin{itemize}[leftmargin=1.2em, itemsep=0.2em]

\item CMPSC 131: Programming and Computation I
\item CMPSC 132: Programming and Computation II
\item CMPSC 221: Object Oriented Programming
\item CMPSC 311: Introduction to Systems Programming
\item CMPSC 360: Discrete Mathematics
\item CMPSC 461: Programming Language Concepts
\item CMPSC 464: Theory of Computation
\item CMPSC 465: Data Structures and Algorithms

\end{itemize}

\vspace{0.4em}

\textbf{Computer Engineering}

\begin{itemize}[leftmargin=1.2em, itemsep=0.2em]

\item CMPEN 270: Digital Design
\item CMPEN 331: Computer Organization and Design

\end{itemize}

%====================================================
% Awards and Honors
%====================================================

\section*{Awards and Honors}

1st Place, ThinkNeuro LLC Project Competition \hfill Apr 2026

1st Place, IEEE Student Chapter Battle Bots Competition \hfill Apr 2026

PIT-UN Funding Award (\$3,750) \hfill Apr 2025

3rd Place, ACM Pathfinder Robotics Challenge \hfill Apr 2025

Republic of Korea Air Force Commendation Award \hfill Aug 2023

%====================================================
% Technical Skills
%====================================================

\section*{Technical Skills}

\textbf{Programming Languages:} Python, C, C++, Java, JavaScript, TypeScript, SQL

\vspace{0.3em}

\textbf{Frameworks and Tools:} React, Node.js, NestJS, Express, Flask, MongoDB, Redis, Firebase, Kubernetes, Git, AWS Bedrock, DuckDB

\vspace{0.3em}

\textbf{Machine Learning:} TensorFlow, PyTorch, FAISS, Mem0, SHAP, spaCy, Scikit-learn, Flower

\end{document}
